\documentclass[11pt]{article}
\usepackage[margin=1in]{geometry}
\usepackage{amsmath,amssymb,mathtools}
\usepackage{fontspec,unicode-math}
\setmainfont{Latin Modern Roman}
\setsansfont{Latin Modern Sans}
\setmathfont{Latin Modern Math}
\usepackage{xurl,hyperref}
\hypersetup{colorlinks=true,urlcolor=blue,linkcolor=blue}
\newcommand{\sref}[2]{\hyperlink{source:#1:#2}{[S#2]}}
\newcommand{\cataloguescope}[1]{\paragraph{Recorded mathematical question.}#1}
\begin{document}
\section{Infinitely many clusters at half density on transient disk triangulations}\label{statement}

\subsection{Definitions and mathematical statement}

\paragraph{Shared notation.}$\mathbb N=\{1,2,\ldots\}$, and $\mathbb Z,\mathbb Q,\mathbb R,\mathbb C$ denote the integers, rationals, reals and complex numbers. Any other conventions are specified locally.


Let $G$ be the connected one-skeleton of a locally finite triangulation of an open topological disc, with a uniform bound on vertex degrees. In Bernoulli site percolation, vertices are independently open with probability $p$, and clusters are the connected components induced by the open vertices. Write $\mathbb P_p$ for this product law and
\[p_c(G)=\inf\{p:\mathbb P_p(\text{some open cluster is infinite})>0\}.\]
Simple random walk moves from its current vertex to a uniformly chosen neighbor. Recurrence means return to its starting vertex infinitely often almost surely; transience means finitely many returns almost surely.
\paragraph{Statement recorded in the supplied source.}
If $G$ is transient, prove that the number of infinite open clusters under $\mathbb P_{1/2}$ is infinite almost surely. \sref{statement}{2}
\subsection{Short English statement}
Must a transient bounded-degree disc triangulation have infinitely many infinite clusters at half density?
\subsection{Sources}
\begin{description}
\item[\hypertarget{source:statement:1}{S1}] ULAM.AI and the UnsolvedMath Contributors, \emph{UnsolvedMath: A Curated Collection of Open Mathematics Problems}, catalogue entry. CC BY 4.0. \url{https://huggingface.co/datasets/ulamai/UnsolvedMath}.
\item[\hypertarget{source:statement:2}{S2}] Itai Benjamini, \emph{Percolation and Coarse Conformal Uniformization}, June 2015; original author manuscript preserved in Arquivo.pt capture of 31 December 2020, Section 2, first paragraph of p. 4. \url{https://arquivo.pt/noFrame/replay/20201231041529id_/http://www.wisdom.weizmann.ac.il/~itai/resistance_percolation.pdf}.
\end{description}
\label{end:statement}
\end{document}
